\subsection{Alternatif Arsitektur Integrasi PGHD - DecMed}

Apa yang bisa dibandingkan:

\begin{enumerate}
    \item Integrasi PGHD → EHR langsung (server-centric)
    \begin{itemize}
        \item Studi AHRQ / JAMIA tentang integrasi PGHD ke EHR; PGHD biasanya dikumpulkan di server aplikasi khusus lalu di-push ke EHR.
        
        \item Kelebihan: sederhana, selaras EHR klasik.
        
        \item Kekurangan: tidak memanfaatkan DLT, tidak meng-address kebutuhan provenance desentral.
    \end{itemize}

    \item PGHD via FHIR server (SMART-on-FHIR, SMART Markers)
    \begin{itemize}
        \item SMART Markers: mengumpulkan PGHD melalui app lalu mengirim ke FHIR server yang terhubung ke EHR.
        
        \item Kelebihan: interoperabilitas tinggi, standar FHIR.
        
        \item Kekurangan: FHIR server biasanya infrastruktur terpusat; bukan DLT.
    \end{itemize}

    \item PGHD di atas blockchain / marketplace PGHD
    \begin{itemize}
        \item Contoh: decentralized marketplace for PGHD (public blockchain, fokus insentif).
        
        \item Kelebihan: provenance \& integritas kuat.
        
        \item Kekurangan: tidak fokus ke integrasi EMR rumah sakit tunggal seperti DecMed, overhead transaksi tinggi.
    \end{itemize}

    \item Pendekatan “PGHD sebagai ekstensi DecMed” (yang kamu pilih)
    \begin{itemize}
        \item Menggunakan IOTA \& IPFS yang sudah menjadi fondasi DecMed, PGHD menjadi tipe data baru yang mengalir ke storage yang sama (IPFS–IOTA), menggunakan CapBAC/PRE yang sudah ada untuk kontrol akses.
        
        \item Di sini kamu tunjukkan bahwa pilihan arsitekturmu bukan random, tapi: memanfaatkan kekuatan DecMed (DLT, IPFS, CapBAC), sambil mengambil praktik baik dari integrasi PGHD–EHR (FHIR, batching, edge processing).
    \end{itemize}
\end{enumerate}

